% =====================================================================
%  Phase C --- The Cloud: store, timestamps, dashboard
% =====================================================================
\section{The Cloud}
\label{sec:phaseC_cloud}

\subsection{One process, two faces}

The Cloud is deliberately a single process presenting two interfaces: an
MQTT client that talks to robots, and an HTTP server that talks to the
operator. A third interface, a console on the main thread, is where
orders are actually typed --- the operator's seat is in the process that
holds the private key.

\paragraph{No web framework} The API has five endpoints and
\texttt{http.server} is in the standard library. A framework would add a
dependency, a version to pin and a deployment story, and would not make
any of the five shorter. The one thing it would buy --- production-grade
concurrency --- is not needed by a dashboard with one operator, and
pretending otherwise would be the kind of choice that looks like
engineering and is not.

\subsection{Storage: JSON Lines, not a database}

The brief asks for storage in a file of appropriate format. Appropriate
here means openable by a human, appendable without rewriting, and
diffable --- which is JSON Lines: one visit per line, append-only, plus
the images written out as real files, plus a CSV rebuilt on demand.

\begin{lstlisting}[style=bash]
store/
  measurements.jsonl        append-only, one visit per line
  photos/P2,5R/2026....jpg  the photograph, as sent
  qr/P2,5R/2026....png      the QR code, as sent
  measurements.csv          rebuilt on export
  requests.csv              one row per order
\end{lstlisting}

A database would be defensible and is what a production system would
use. It would also mean the evidence for every number in this report
lives inside a binary file nobody can inspect during a demonstration.
For a system whose entire claim is that the data arrived intact, being
able to open the file and read the line is worth more than query speed
over a few thousand rows.

The images come back out of the JSON because leaving them there would
make the file unreadable --- a 6~kB blob per line --- and would mean the
only way to look at a photograph is to write a program. Extracted, the
line stays about 400 bytes and legible.

One corrupt line does not cost the campaign: the loader names the line
it skipped and keeps the rest, because a truncated final line is what a
kill during a write looks like and it is recoverable.

\subsection{Three timestamps, because there are three moments}

A reading has three distinct times, and the system originally kept one.
The only honest answer to ``how long did that sweep take'' was then to
watch a clock while it ran.

\begin{table}[!t]
\caption{The three moments a stored reading carries, each on the clock
of the machine that owns it. End-to-end latency is derivable from the
stored line alone.}
\label{tab:timestamps}
\centering\footnotesize
\begin{tabular}{@{}ll@{}}
\toprule
field & the moment it records \\
\midrule
\param{request\_issued\_at} & the Cloud signed the order \\
\param{measured\_at}        & the robot read the plant \\
\param{received\_at}        & the Cloud filed the report \\
\param{latency\_s}          & derived: issued $\rightarrow$ received \\
\bottomrule
\end{tabular}
\end{table}

The column \param{timestamp} was removed rather than redefined, and the
suite asserts it is gone --- with three moments in play, a field called
``timestamp'' is a question, not an answer.

Orders are tracked in parallel in \filename{requests.csv}, with
\param{issued\_at}, \param{completed\_at} and \param{elapsed\_s}. A
mission that ends \emph{without} filing every report still ends: the
record closes on the terminal acknowledgement, so an abandoned sweep
reports how long it ran before giving up rather than showing as running
forever. The dashboard distinguishes the two verbs --- ``done'' versus
``gave up'' --- because the same numbers under the same word would
report a failure as a success in the one place an operator looks for the
answer.

\subsection{The dashboard}

The dashboard is a single page with no build step and no framework,
polling \texttt{/api/state} every two seconds. It shows the greenhouse
as a scale map with all 48 stations, the readings of the selected
quantity, the plant photograph and QR image, live request progress, and
the rejection panel.

Two problems in it are worth recording because both are geometric rather
than aesthetic:

\begin{itemize}[leftmargin=*,itemsep=1pt,topsep=2pt]
  \item \textbf{The two stations of an inner aisle collide when drawn.}
        At \SI{0.10}{m} apart their marks touch at map scale. Each value
        chip is pushed away from its own row and joined to its mark by a
        stem, and the suite asserts the plain rendering really does
        collide and the corrected one does not --- so the fix cannot be
        removed by someone who thinks it is decoration.
  \item \textbf{Text inherits the map's flip.} The SVG applies a
        $y$-negating transform so that $+y$ is north; every text element
        must counter-flip it or read mirrored. That counter-flip is
        written once, in a helper, and the suite asserts it is not
        open-coded per call site.
\end{itemize}

\paragraph{Access control} The dashboard API is protected by a random
token printed at startup. \texttt{/media/} is behind the same token:
photographs, QR images and the CSV export are measurements, so they sit
behind the same gate as the API that describes them.

The token is offered once, in the URL, and the server's first act is to
get rid of it: it sets an \texttt{HttpOnly; SameSite=Strict} cookie and
redirects to the same page with the query string stripped. A credential
that authorises driving the robot has no business living in the address
bar, the browser history, the \texttt{Referer} of every outbound link
and any proxy log in between. Repeated wrong tokens from one address
earn a bounded lockout.

Authentication can still be disabled with
\texttt{-{}-dashboard-token ''}, but only together with
\texttt{-{}-http-host 127.0.0.1}; anything else refuses to start. It
used to print a warning and serve anyway, and a warning at startup is
read once and scrolls away. Sec.~\ref{sec:security_posture} covers this
and the rest of the posture.

\subsection{Stopping cleanly}

The operator can stop the session from the console with Ctrl-C or from
the browser with \textsc{quitter}. Both paths run \emph{one} teardown
function behind a lock, because both can happen at once and exporting
the CSV twice from two threads truncates the file the first one is
writing.

The web path additionally has to exit a process whose main thread is
blocked inside \texttt{input()}, which no exception raised on an HTTP
worker thread can interrupt. It replies first, flushes the socket, and
then terminates from a short-delayed thread --- exiting from inside the
handler resets the connection often enough to matter, and the browser
then reports a failure for a shutdown that in fact worked.

The offline demo serves the same dashboard with the shutdown hook set to
\texttt{None}, and the suite asserts that it refuses the request rather
than killing its caller's interpreter.
